\chapter{{Evaluative} {meaning} {of} {postnominal} {\textit{cualquiera} }{in} {European} {Spanish}}\label{sec:3}

As I already showed in \chapref{sec:2}, the evaluative reading of (\textsc{un)} N \textit{cualquiera} ‘common, ordinary’ can appear after predicative verbs, such as the copular(-like) verbs \textit{ser} ‘to be’ and \textit{parecer} ‘to seem’ (see \sectref{sec:2.1.2}). In this chapter, I investigate the evaluative meaning of (\textsc{un)} N \textit{cualquiera} in more detail. In order to do so, I will start with the description of (\textsc{un)} N \textit{cualquiera} in predicative contexts and extend this description to other verbs.

{My observations are primarily based on corpus data from \textit{Corpus del Español} (CdE), in particular the following subcorpus CdE\_web dialects\_EurSp. On particular occasions, I also use speaker judgments.} The crucial result of my observations, which will serve the analysis of \textit{cualquiera}, is that predicative \textsc{un} N \textit{cualquiera} appears extremely often (ca. 98\% of all data), especially under negation (around 80\%), see \sectref{sec:3.1}.

As I will argue in the subsequent sections (\sectref{sec:3.2}–\sectref{sec:3.5}), negation and the predicative position of \textsc{un} N \textit{cualquiera} trigger the reading ‘a common N’ or ‘an ordinary N’ which I subsumed under the term EVAL to distinguish it from modal readings such as ‘any’ and ‘random’. In order to explain certain types of data, such as coordination with adjectives, I will postulate a reanalysis of \textit{cualquiera} into a (frequency) adjective (i.e. with the meaning ‘common’) in European Spanish (see \sectref{sec:3.5}). I will argue that negation plays an important role in this reanalysis, as it introduces a contrastive property into the discourse, which strengthens the property interpretation of postnominal \textit{cualquiera} as ‘common’ or ‘unremarkable’.

The chapter is organised as follows. In \sectref{sec:3.1}, I first present some data concerning the distribution of postnominal \textit{cualquiera}, which show that it typically occurs in predicative syntactic contexts, and, most frequently, under negation. I will then describe the evaluative function in more detail in \sectref{sec:3.2}. I will argue that what I call the ‘common’ reading can be assumed as the basic meaning of postnominal \textit{cualquiera}, from which the derogatory meaning of ‘low value’ can be derived. In \sectref{sec:3.3}, I present and discuss some additional properties of postnominal \textit{cualquiera}, in particular the possibility of an overt experiencer, the extensional meaning of the ‘common’ reading and its relationship to the free choice reading, and the possibility of postnominal \textit{cualquiera} under modal verbs. After a brief intermediate summary in \sectref{sec:3.4}, I will show in \sectref{sec:3.5} that \textit{cualquiera} in its ‘common’ reading can be derived from the FCI function if we analyse it as a quantifier under a generic operator in predicative and episodic contexts. \sectref{sec:3.3} will offer an analysis according to which \textit{cualquiera}, when it occurs under definite nouns, is no longer interpreted as a quantifier, but as an item referring to a property of types of individuals. I will argue that both analyses, the quantifier analysis and the adjectival analysis, are needed in order to capture the corpus data and the ambiguity between the free choice and the evaluative interpretation of \textit{cualquiera} observed in the literature (see \chapref{sec:2}, in particular my discussion of \citealt{Alonso-OvalleMenéndez-Benito2011,am2018}).

\section[{Distribution} {in} {the} {Corpus} {del} {Espanol\_web} {dialects}]{{Distribution} {of} {postnominal} {\textit{cualquiera} }{in} {the} {Corpus} {del} {Espanol\_web} {dialects}} \label{sec:3.1}

In this subsection, I only used the subcorpus Corpus del Espanol\_web dialects (henceforth: CdE\_web dialects) because other subcorpora did not allow for the automatic extraction of frequencies in European Spanish by the time I collected data in CdE in 2016.\footnote{{In the meantime, many larger subcorpora have emerged that also allow for automatic frequency extraction. In the future, I will make a comparison between these subcorpora.}} Henceforth, I will use square brackets [ ] for the different constructions I searched for in this corpus, e.g. [V \textsc{un} N \textit{cualquiera}] (= postnominal \textit{cualquiera} after the verb).

\tabref{tab:key:6} summarises the distribution of +/– affirmative mood of [V \textsc{un} N \textit{cualquiera}]. It appears in negative contexts (with implicit or explicit negation) more often than in affirmative ones. 


\begin{table}
\begin{tabularx}{\textwidth}{rYY}
\lsptoprule
 {V} {+} {\textsc{un}} {N} {\textit{cualquiera}} & {V} {+} {\textsc{un}} {N} {\textit{cualquiera}} {in} {affirmative} {contexts} & {Neg} {+} {V} {+} {\textsc{un}} {N} {\textit{cualquiera}}\\
 \midrule
 100\% & 18\% & 82\%\\
 353 & 66 & { 287}\newline { (explicit negation = 283 = 99\%}  of total negation)\\
\lspbottomrule
\end{tabularx}
\caption{\label{tab:key:6} Distribution of +/– affirmative contexts of V+ \textsc{un} N \textit{cualquiera} in CdE\_web dialects in European Spanish. The first line shows the percentages and the second line the absolute numbers of occurrences.}
\end{table}

The most frequent verb type that precedes \textsc{un} N \textit{cualquiera} is the copular verb \textit{ser} ‘be’ (98\% of all verbs), whose most frequent verb form (verb token) is the indicative present tense (\textit{es} ‘is’, 88\%). The less frequent verb type are transitive verbs such as \textit{poner} ‘put’ (\tabref{tab:key:7}).

\begin{table}
\begin{tabularx}{\textwidth}{rYY}
\lsptoprule
 {V} {+} {\textsc{un}} {N} {\textit{cualquiera}} & {Ser} {+} {\textsc{un}} {N} {\textit{cualquiera}} & {Other} {verbs}\\
 \midrule
 100\% & 98\% & 2\%\\
 353 & 345 & 8\\
& \textit{es} (88\% = 308) & \textit{puse, prestó}\\
\lspbottomrule
\end{tabularx}
\caption{\label{tab:key:7} Distribution of verbs in V + \textsc{un} N \textit{cualquiera} in CdE\_web dialects in European Spanish. Other determiners were disregarded due to low frequency. The first line shows the percentages and the second line the absolute numbers of occurrences.}
\end{table}

The observation that the most frequent verb token is \textit{es} 3\textsuperscript{rd} person singular and not plural suggests that \textsc{un} N \textit{cualquiera} is usually used by speakers to describe a property of a single person or entity and not a property of different entities or people. As we will see in \sectref{sec:3.2}, the most frequent token is \textit{un día cualquiera} ‘an ordinary day’. Thus, \textsc{un} N \textit{cualquiera} is mostly used to describe the day of speaker’s reference as ordinary or usual.

All examples that I have examined show that 82\% of [V \textsc{un} N \textit{cualquiera}] occur under implicit negation as in \REF{bkm:Ref42868256} \footnote{{I counted as implicit negation only those cases where the context was very clear that the speaker does not believe that the proposition expressed by V} \textsc{un} {N} {\textit{cualquiera}} {is true.}} or explicit negation as in \REF{bkm:Ref42868258}, whereby explicit negation is a more frequent type (99\% of total cases with negation):

\largerpage[2]
\ea \label{bkm:Ref42868256} %ex 122 
\gll ¿Sigues pensando que fue un día \textit{cualquiera}? […] es lunes 5 de diciembre de 2011.\\
do.still thinking that was a day \textsc{cualquiera} […] is Monday 5 of December of 2011\\
\glt ‘Are you still thinking that it was a normal day? […] it’s Monday, December 5, 2011.’

  (@ 40)
\z

\ea \label{bkm:Ref42868258} %ex 123
\gll […] que esto no es un libro \textit{cualquiera}, esto es un libro de James Dashner.\hspace*{-4mm}\\
[…] that this not is a book \textsc{cualquiera} this is a book of James Dashner\hspace*{-4mm}\\
\glt ‘[…] that this is not an ordinary book, this is a book from James Dashner.’

  (@ 41)
\z\clearpage

Affirmative contexts appear less frequently than negative contexts (see \tabref{tab:key:5}). Recall that affirmative contexts usually represent a puzzle for the licensing of \textit{cualquiera} as an FCI (see \sectref{sec:2.2}) unless rescued by some operators (see \sectref{sec:2.3}).

In affirmative contexts (i.e. contexts without explicit or implicit negation) the verb in [V \textsc{un} N \textit{cualquiera}] can be of different types, such as transitive verbs and predicative verbs, that is, verbs that take predicative complements as their argument. Predicative verbs are the most common verb type in this corpus as will be shown in \tabref{tab:key:6}. They are often realised as a copular verb in present indicative, as in \REF{bkm:Ref42868883} and \REF{bkm:Ref42868888}:

\ea \label{bkm:Ref42868883} %ex 124
\gll Parece una chiquilla \textit{cualquiera}, vestida con una túnica liviana que adorna su cuerpo hermoso y esconde el brazo derecho mutilado.\\
seems a girl \textsc{cualquiera} dressed with a tunic lightweight that adorns her body beautiful and hides the arm right mutilated\\
\glt ‘She looks like a common little girl, ...’

  (@ 42)
\z


\ea \label{bkm:Ref42868888} %ex 125
\gll De todos modos, tampoco os esperéis la gran historia. Es una historia \textit{cualquiera}.\\
of all ways neither expect you the great story is a story \textsc{cualquiera}\\
\glt ‘Anyways, don’t expect a big story either. It’s a story like any other.’

  (@ 43)
\z

There are fewer occurrences of other verb types that are not copular verbs, e.g. transitive verbs such as \textit{prestó} and \textit{puse} in \REF{bkm:Ref42868870} and \REF{bkm:Ref42869017}. Non-copular verbs that appear in [V \textsc{un} N \textit{cualquiera}] are usually realised with perfective aspect, expressed in Spanish with the \textit{pretérito indefinido} (in so-called episodic contexts; see section \sectref{sec:2.3} for discussion of episodic contexts):

\ea \label{bkm:Ref42868870} %ex 126
\gll La NYK \textit{prestó} un contenedor \textit{cualquiera} a la BBC.\\
the NYK lent a container \textsc{cualquiera} to the BBC\\
\glt ‘The NYK borrowed a random container to the BBC.’

  (@ 44)
\z


\ea \label{bkm:Ref42869017} %ex 127  
\gll Le \textit{puse} un nombre \textit{cualquiera}.\\
I gave a name \textsc{cualquiera}\\
\glt ‘I put some random name.’

  (@ 45)
\z

{With respect to the distribution of the determiner preceding N \textit{cualquiera}, the most frequent determiner in the corpus under investigation is the indefinite determiner \textsc{un}, as already observed in the literature (see \chapref{sec:2}). However, there are few cases of N \textit{cualquiera} being preceded by definite determiners, demonstratives, and quantifiers (definite determiner N \textit{cualquiera} 136 = 0.30 per mil vs. indefinite N \textit{cualquiera} 1541 = 3.36 per mil).}

In the following examples, the definite noun refers to a certain kind of man, namely to the common man {as a type}. One could also describe this type as \textit{el hombre de la calle}, lit. ‘the man from the street’ or ‘common/normal man’:

\ea \label{ex:key:150}
\gll  El hombre \textit{cualquiera} del autobús no siente vergüenza alguna, aunque debería,…\\
the man \textsc{cualquiera} of.the bus \textsc{neg} feels shame any although should\\
\glt ‘[…] The common man from the bus does not feel any shame, although he should,…’

  (@ 82)
  \z

\ea \label{ex:key:189}
\gll Las armas de este hombre \textit{cualquiera}, de aspecto normal, se resumen en: prudencia.\\
the weapons of this man \textsc{cualquiera} of appearance normal \textsc{refl} summarised in: prudence\\
\glt ‘The weapons of this common, normal-looking man are summarised in: caution.’

  (@ 83)
\z

\textit{Cualquiera} also occurs with quantified nouns (see \sectref{sec:2.1.3} on the discussion concerning quantified nouns):

\ea \label{bkm:Ref58144031} %ex 190
\gll Esto lo saben mucha gente \textit{cualquiera}.\\
this it know many people \textsc{cualquiera}\\
\glt ‘This is known to many ordinary people.’

  (@ 84)
\z


\ea \label{bkm:Ref63546096} %ex 191
\gll Las cifras del IMC se obtienen tras el estudio de los datos de muchísimas chicas \textit{cualquiera}.\\
the figures of.the BMI \textsc{refl} obtener.\textsc{3pl.prs} after the study of the data of very.many girls \textsc{cualquiera}\\
\glt ‘The figures in the IMC are obtained after studying many ordinary young women.’

  (Fulton, \textit{Seis semanas}, El índice de masa corporal)
\z 

\ea \label{ex:key:192}
\gll vamos a morir algún día \textit{cualquiera}\\
go.\textsc{1PL} to die some day \textsc{cualquiera}\\
\glt ‘we will die on some ordinary day’

  (@ 85, @ 86)
\z
 

It has been shown that postnominal \textit{cualquiera} occurs more often under negation than in affirmative contexts (see \tabref{tab:key:6}). The question arises as to what the interpretation of postnominal \textit{cualquiera} is in affirmative and negative contexts. I will analyse the interpretations in these contexts in more detail in \sectref{sec:3.2}.

In the following sections, I will first focus on \textsc{un} N \textit{cualquiera} as opposed to other determiners, as it is the most frequent type in European Spanish.  As predicative \textsc{un} N \textit{cualquiera} is a very important type in this variety according to my corpus analysis, I will investigate this type more closely, especially the meanings of \textsc{un} N \textit{cualquiera} under predicative verbs. I will then provide an analysis of \textit{cualquiera} under definite nouns. 

\section{{Interpretations} {of} {postnominal} {\textit{cualquiera}}} \label{sec:3.2}

{I have analysed all examples of postnominal \textit{cualquiera} and will present the results by focusing on the interpretations of the most frequent tokens after the verb in affirmative and negative contexts. One heuristic for approaching the semantic readings of \textsc{un} N \textit{cualquiera} is to investigate the range and types of appositives and adverbial modifications that co-occur with \textsc{un} N \textit{cualquiera}, as they can provide us with additional information about the syntactic and semantic status of \textit{cualquiera} (see also \citealt{Condoravdi2015}: 227 for English). I have looked at the possible appositives and adverbial modifications that \textsc{un} \textit{N cualquiera} can have within the corpus CdE\_web dialects as well as the pragmatic context in which this construction appears. I do not report the frequency of appositives or adverbial modification as there are only few examples (1 or 2) of a particular type of appositives or adverbial modification. However, I do report the frequency of interpretations that I identified via adverbial modification, appositives, and pragmatic context.}

%I will show that \textsc{un} N \textit{cualquiera} has mostly the evaluative interpretation (EVAL), which has different flavours depending on the context (COM(mon), DER(ogatory), SIM(ilarity)). As will be shown in the following section, the most frequent flavour of EVAL is SIM. This reading has not been observed previously in the literature, in contrast to the ‘unremarkable’ interpretation of \textit{cualquiera} (see \sectref{sec:2.3}). Even the ‘unremarkable’ interpretation, which has been already observed in the literature in relation with \textit{cualquiera} (see \sectref{sec:2.3}), appears in two different flavours: COM and DER. These different flavours of the ‘unremarkable’ interpretation have not been studied before in the literature in relation with \textit{cualquiera}. My study of these different flavours of EVAL is, to my knowledge, the first study to investigate EVAL on the basis of corpus data.

\subsection{{Readings} {in} {affirmative} {contexts}} \label{sec:3.2.1}

{There is a total of 187 occurrences of postnominal \textit{cualquiera} in affirmative contexts with the search algorithm [!no Verb Article Noun cualquiera]. The ten most frequent instances in the corpus are: \textit {es un día cualquiera, es un chico cualquiera, era un día cualquiera, acaban comprando un ereader cualquiera, narra un día cualquiera, para coger un libro cualquiera, dado un número cualquiera, empecé una tarde cualquiera, coges una foto cualquiera, ha un servicio cualquiera} (English {‘it's a normal day, it's a normal boy, it was a normal day, they end up buying some (random) ereader, it's a normal day, to take any book, given a random number, a normal afternoon began, you take any/random photo, there is a normal service.’})}

As can be seen from the ten most frequent occurrences in affirmative contexts, postnominal \textit{cualquiera} appears under indicative verbal mood, e.g. \textit{es/era/empecé} ‘is/was/began …’ or under non-indicative verbal mood, e.g. infinitives, imperatives, future, etc. as in \textit{para coger, coges} ‘to take, you choose’. The present form \textit{coges} has an imperative interpretation in the example \REF{ex:key:Coges}. Under non-indicative verbal mood, the interpretation is the FCI ‘any’ or RCI ‘random’:

\ea \label{ex:key:Leéis} %ex 
    
     \gll ¿Leéis el prólogo antes que el libro o lo dejáis para el final? \textbf{Hoy} \textbf{es} \textbf{un} \textbf{gran} \textbf{día} \textbf{para} \textbf{coger} \textbf{un} \textbf{libro} \textit{\textbf{cualquiera}} y sumergirse en sus páginas.\\
     \textbf{read.PRS.2PL} the prologue before that the book or it \textbf{leave.PRS.2PL} for the end today is a great day to take a book \textsc{cualquiera} and immerse in its pages\\
%today is a good day to buy a book, to participate in the signing of books or to approach \textsc{refl} to a library\\
    \glt ‘Do you (plural) read the prologue before the book or do you leave it for the end? Today is a big day to take any/a random book...' (@ 286)  
\z

\ea \label{ex:key:Coges} %ex 
    \gll \textbf{Coges} \textbf{una} \textbf{foto} \textit{\textbf{cualquiera}}, de \textit{cualquier} cosa, le metes un par de filtros y ya tienes una obra de arte.\\
    %\textsc{cualquiera} book\\
    take a photo \textsc{cualquiera} of \textsc{cualquier} thing, it.\textsc{dat} put a pair of filters and already have a work of art\\
    \glt ‘Take any foto/a random foto, of anything, […]’.
    
    (@ 286) 
\z

{Under indicative verbal mood, postnominal \textit{cualquiera} has the interpretation of ‘normal’, ‘common’, ‘ordinary’ or ‘typical’ (henceforth COM) in all examples except three, where it has the derogatary interpretation of ‘low value’ (henceforth DER) or the random interpretation (RCI):} 

\ea \label{ex:key:Aaron} %ex 
  \gll Aaron parece un hombre \textit{cualquiera}, acomodado a la rutina del barrio donde vive, ocupado en su trabajo.\\
    Aaron seems a man \textsc{cualquiera}, adjusted to the routine of.the neighborhood where lives busy in his work\\
    %\textsc{cualquiera} book\\
    \glt ‘Aaron looks like a normal man, adjusted to the routine of his neighborhood, immersed in his work’. (@ 279) 
\z

\ea \label{ex:key:Hola} %ex 
      - Hola, soy Celia. - ¡¡Hola Celia!! - Tengo 23 años y soy adicta. - ¡¡Bienvenida!! ¡Cuéntanos tu historia! - 
     
     \gll Empecé una tarde \textit{cualquiera} de un mes\textit{cualquiera}.\\ 
     started a afternoon \textsc{cualquiera} of a month \textsc{cualquiera}.\\
     
    La llamada de La lista era fuerte y poderosa.
    
    \glt ‘Hello, I'm Celia. Hello Celia. I'm 23 years old and I'm an addict. Welcome. Tell us your story. I started on an ordinary day of an ordinary month.' 
    (@ 283)  
\z	

{Two examples under indicative verbal mood have an explicit derogatory interpretation, which is identifiable by expressions like ‘there are worse days’ or ‘mediocre’.}

\ea \label{ex:key:En } %ex 
       En fin, baños y cenas (otra guerra, ya me duché ayer, no me laves la cabeza, la hermana me ha escupido…..). A las nueve más o menos a la cama, los dientes y los cuentos; mientras ella les lee los cuentos yo recojo por tercera o cuarta vez la cocina; cuando acabo les voy a hacer unos arrumacos a los niños en la cama, les gusta que me despida de ellos. Terminado; al sofá, soma, quiero soma, y enciendo la puta tele, me alieno un poco y si no caigo rendido, me paso al ordenador, a trabajar un poco en mis cosas, o a desenchufar me un poco más leyendo blogs como el tuyo.\\
     \gll \textbf{Este} \textbf{es} \textbf{un} \textbf{día} \textbf{\textit{cualquiera}}, \textbf{pero} \textbf{te} \textbf{juro} \textbf{que} \textbf{los} \textbf{hay} \textbf{mucho} \textbf{peores}.\\
     this is a day \textsc{cualquiera} but to.you swear that them are much worse\\
    \glt ‘This is an ordinary day, but I swear to you that there are much worse days.' (@ 278)  
\z

\ea \label{ex:key:Por}
      Por todo esto se podría decir que \textbf{es una comedia \textit{cualquiera}} de las que nos han llegado en los últimos tiempos, y realmente eso es lo que es. La diferencia es que a esta no le sobran minutos ni escenas, que el ritmo no baja demasiado y no llega a aburrir\_te, y, sobre todo, que hasta consigue robar te alguna carcajada, lo que cada vez es más difícil con las comedias de Hollywood.\\
     \gll Solo por eso, \textbf{una} \textbf{película} \textbf{mediocre} \textbf{como} \textbf{esta} ya merece más la pena que muchas otras.\\
only for that a movie mediocre like this already deserves more the worth than many others\\
    \glt ‘For all these reasons we can say that \textbf{it's an ordinary comedy} among those that we had in recent time […] Only for this reason, \textbf{a mediocre movie like this one} is more worth than other movies.] (@ 284)’
\z

{In only one example, \textit{cualquiera} has the interpretation of ‘random’ under the transitive verb ‘buy’ (see §2 for RCI):}

\ea \label{ex:key:mis} %ex 
    \gll mis consultantes acaban comprando un ereader \textit{cualquiera} (el {más barato)} y cargándolo de libros en su ordenador, [...]\\
    my clients end.up buying an ereader \textsc{cualquiera} (the {cheapest)} and {loading it} with books in their computer [...]\\
    \glt ‘my consultants just bought some random ereader (the cheapest one) and […],' (@ 280)
\z

The affirmative sentence \textit{Es} \textit{un} N\textsubscript{[animate]} \textit{cualquiera} can appear with a comparative expression which expresses similarity between the referent to which the property \textsc{un} N\textsubscript{[animate]} \textit{cualquiera} applies (here: \textit{Juan}) and other individuals of the same comparison class (here: ‘everyone else like you and me’):

\ea \label{ex:key:128} %ex 128
\gll Juan es un hombre \textit{cualquiera}, cargado de errores. Como tú, o como yo.\\
Juan is a man \textsc{cualquiera} loaded with errors like you, or like I\\
\glt ‘Juan is a human being like everyone else, full of flaws. Like you, or like me.’

  (@ 46)
\z

The similarity construction is expressed by the comparison term \textit{como} ‘like’, which expresses similarity between the referent and a high number of individuals that share some (basic or prototypical)\footnote{{The assumption that} {\textit{common}} {properties must be basic follows from the similarity of all individuals denoted by N} {\textit{cualquiera}}{. For instance,} {\textit{un trabajo cualquiera ‘a job \textsc{cualquiera}’}} {is a job that shares with other jobs the most basic properties, whatever it is that defines a job as a job, such as getting up early, working for eight hours, going back home, etc.}} properties or at least one prototypical property with the subject. It is impossible to find a comparison to a set of individuals that contains only one individual or none:

\ea \label{ex:key:129} %ex 129
\gll Juan\textsubscript{j} es un hombre \textit{cualquiera}, cargado de errores. \# No hay otro hombre que tenga errores como él\textsubscript{j}.\\
Juan is a man \textsc{cualquiera} loaded with errors. \# \textsc{neg} there.is another man who has errors like he\\
\glt ‘Juan is a human being like everyone else, full of flaws. \#There is no one else like him.’
\z

In the following example, the morning the speaker refers to is a (kind of) morning that is similar to all or many other mornings.

\ea \label{bkm:Ref42870625} %ex 130
\gll Fue una mañana \textit{cualquiera}, como cada mañana. Había estado leyendo y llamó uno de los muchachos de la cantina.\\
was a morning \textsc{cualquiera} like every morning had been reading and called one of the boys from the bar\\
\glt ‘It was a morning like any other mornings. I had been reading and one of the guys from the canteen called.’

  (@ 46)
\z

Indeed, native speakers deem it possible for \REF{bkm:Ref42870616} to have the same interpretation as in \REF{bkm:Ref42870625}. This shows that even though the comparison is not overtly or explicitly expressed by \textit{como} ‘like’, speakers can infer an implicit comparison from \textsc{un} N \textit{cualquiera} (see \sectref{sec:3.3.3}): 

\ea \label{bkm:Ref42870616} %ex 131 
\gll Fue una mañana \textit{cualquiera}. Había estado leyendo y llamó uno de los muchachos de la cantina.\\
was a morning \textsc{cualquiera} had been reading and called one of the boys from the bar\\
\glt ‘It was a morning like any other. I had been reading and one of the guys from the canteen called.’

  (@ 47)
\z

There are also appositives that describe N \textit{cualquiera} as N \textit{corriente} ‘normal/regular/frequent N’ as in \REF{bkm:Ref57554115} or N \textit{cualquiera} itself can appear as an appositive clause that describes N \textit{normal} ‘normal N’ as in \REF{bkm:Ref63545479}. 
%\footnote{{The distinction between SIM and COM is primarily based on different types of appositives and paraphrases, i.e. SIM refers to similarity constructions like \textit{como cualquiera} that contain an overt quantifier, whereas COM refers to appositive adjectives like \textit{corriente}. However, semantically SIM and COM are related as argued above because 'being similar like anyone/anything else' implies 'being common'}}

\ea \label{bkm:Ref57554115} %ex 132 
\gll Hoy es un día \textit{cualquiera}, un día corriente, […]\\
today is a day \textsc{cualquiera} a day common\\
\glt ‘Today is a day like many other days, a regular day, …’

  (@ 48)
\z


\ea \label{bkm:Ref63545479} %ex 133 
\gll Yo soy una persona normal, un tipo \textit{cualquiera}.\\
I am a person normal a guy \textsc{cualquiera}\\
\glt ‘I am a normal person, a guy like other guys.’

  (@ 49, @ 50)
\z

Finally, some appositives clearly indicate that \textsc{un} N \textit{cualquiera} can have a derogatory meaning (DER). In these contexts, N \textit{cualquiera} is described as not worth mentioning:

\ea \label{ex:key:134} %ex 134
\gll era un tipo \textit{cualquiera}, sin ninguna particularidad digna de mención.\\
was a guy \textsc{cualquiera} without any particularity worthy of mention\\
\glt ‘he was an ordinary guy, without any particular quality worth mentioning.’

  (@ 51)
\z

In the following example, the derogatory reading is expressed by the use of the adverb \textit{simplemente} ‘just’ or \textit{nada más} ‘nothing but’, which expresses that the person described by N \textit{cualquiera} has no properties other than that described by N \textit{cualquiera}:

\ea \label{ex:key:135} %ex 135
\gll Era simplemente un tipo \textit{cualquiera}.\\
was simply a guy \textsc{cualquiera}\\
\glt ‘It was just an ordinary guy.’

  (@ 52)
\z

\ea \label{ex:key:136} %ex 136
\gll No soy nada más que un chico \textit{cualquiera}, que pasaría desapercibido, sin más, normal, que se diría.\\
\textsc{neg} am nothing more than a boy \textsc{cualquiera}, who would.pass unnoticed, without more normal than \textsc{refl} said\\
\glt ‘I am nothing but an ordinary boy, that would go unnoticed, nothing more, normal, that one would say.’

  (@ 53)
\z

The meanings identified in affirmative contexts under indicative verbal mood are related to each other. ‘Being similar to anyone else’ implies ‘to be normal’. The derogatory reading expresses a negative judgment of ‘being normal’ or ‘similar to anyone else’, which can be paraphrased as ‘being just normal’ and ‘not special’. I will investigate this semantic relation more closely in \sectref{sec:3.3.2}.


To summarise the interpretations of \textsc{un} N \textit{cualquiera} in affirmative contexts, \textit{cualquiera} means ‘any’ (FCI) under non-indicative verbal mood. It has the meaning ‘similar like any other’, ‘common’, ‘typical’ (COM) or a derogatory interpretation as ‘not special’, ‘just common’, ‘just like any other’ (EVAL1) under indicative verbal mood under predicative verbs like ‘be’ and ‘appear’. In affirmative contexts, the derogatory interpretation is less common than the neutral interpretation of ‘common’, ‘similar like any other’ (1 \%, i.e. 2 out of 188 examples). In predicative contexts, \textsc{un} N \textit{cualquiera} refers to entities or individuals that share common properties or properties that everyone else has. Thus, \textit{es un día} \textit{cualquiera} expresses ‘it's a day like any other day’ or ‘it's a common day’. A formal analysis of \textit{N cualquiera} as a type/property of individuals will be proposed in \sectref{sec:3.3.1}.

\subsection{{Readings} {under} {explicit} {and} {implicit} {negation}} \label{sec:3.2.2}

{In all examples in the corpus that occur in negative contexts with the search algorithm [no Verb DET N \textit{cualquiera}] (288 examples in total, see \sectref{sec:3.1}), postnominal \textit{N cualquiera} expresses a contrast between \textit{no Verb DET N cualquiera} and an alternative with a property of being special (e.g., \textit{no es un día cualquiera, es un día especial.} ‘It's not a normal day, it's a special day’). 

The first ten most frequent instances of the search algorithm are: {\textit{no es un día cualquiera} ‘not an ordinary day’, \textit{no es un libro cualquiera} ‘not an ordinary book’, \textit{no es un momento cualquiera} ‘not an ordinary moment’, \textit{no es un país cualquiera} ‘not an ordinary country’, \textit{no es un partido cualquiera} ‘not an ordinary party’, \textit{no es una historia cualquiera} ‘not an ordinary story’, \textit{no es un ejército cualquiera} ‘not an ordinary army’, \textit{no es un circo cualquiera} ‘not an ordinary circus’, \textit{no es un jugador cualquiera} ‘not an ordinary player’, \textit{no es un regalo cualquiera} ‘not an ordinary gift’}. 

Under implicit negation, for example with the counterfactual interpretation of the verb \textit{sería} in \REF{bkm:Ref54634548} or the question context in \REF{bkm:Ref54634539}, \textsc{un} N \textit{cualquiera} has an evaluative interpretation:

\ea \label{bkm:Ref54634548}  %ex 137
\gll Pensaron que sería un día \textit{cualquiera}. Pero no. La basura en el mar hizo que ese viernes sea diferente.\\
thought that be.\textsc{cond.3sg} a day \textsc{cualquiera} but \textsc{neg} the trash in the sea made that that Friday be different\\
\glt ‘They thought that it would have been a normal day. But no. The rubbish in the sea turned this Friday into a special day.’

  (@ 54, @ 55)
\z

\ea \label{bkm:Ref54634539} %ex 138
\gll ¿Pensabas que este sábado sería un sábado \textit{cualquiera}?\\
thought that this Saturday be.\textsc{cond.3sg} a Saturday \textsc{cualquiera}\\
\glt ‘You thought that this Saturday would be a normal Saturday?’

  (@ 56, @ 57)
\z

\textit{N cualquiera} under (explicit or implicit) negation is used to express \textit{a contrast}\textbf{ }to a (proper) noun or some other DP that contains a reference to a single individual or type of individual. This individual/entity has the property of being ‘particular’, ‘distinguished’, ‘uncommon’, or having ‘higher value’ than \textit{N cualquiera}, such as Messi (who is known to be an extraordinary football player) in \REF{bkm:Ref66615171} or the author John Dashner\footnote{The further context of \REF{bkm:Ref150779604} (from a book review) reveals that this author is particularly distinguished for writing books full of suspense and action.} in \REF{bkm:Ref150779604}.

\ea \label{bkm:Ref66615171}  %ex 139

\gll Sin Messi, el Barcelona sería un equipo \textit{cualquiera}.\\
without Messi the Barcelona be.\textsc{cond.3sg} a team \textsc{cualquiera}\\
\glt ‘Without Messi, the team Barcelona would be some unimportant football team.’

  (@ 58)
\z

\ea \label{bkm:Ref150779604} %ex 140
\gll […] que esto no es un libro \textit{cualquiera}, esto es un libro de James Dashner.\\
[…] that this \textsc{neg} is a book \textsc{cualquiera} this is a book by James Dashner\\
\glt ‘[…] that this is no ordinary book, this is a book by James Dashner.’

  (@ 41)
\z

To sum up, the negation triggers a contrastive clause expressed by a proper noun (e.g. Messi) or some DP (\textit{un libro de James Dashner}) which singles out one and only one alternative that has the property of being particular, uncommon, or outstanding. This alternative is opposed to the individuals described by \textsc{un} N \textit{cualquiera}.

As I will suggest in my analysis, predicative verbs represent a special linguistic context that gives rise to reanalysis of N \textit{cualquiera} into a property type and thus a shift from a quantifier \textit{cualquiera} into a property \textit{cualquiera} (see \ref{sec:3.3.1}).

\subsection{{The} {semantic/pragmatic} {status} {of} {meanings}} \label{sec:3.2.3}

The question arises as to which of the readings identified in \sectref{sec:3.2.1} and \sectref{sec:3.2.2}, namely the readings ‘any’, ‘common’ or ‘low value’ is a lexical meaning of \textit{cualquiera} and which one is just an implicature and then, whether one can cancel this implicature and thus can cancel the meaning of postnominal \textit{cualquiera} (see \citealt{Aloni2021}, \citealt{am2018} on different views).

{As has been argued by \citet{Grice1975}, conversational implicatures are usually cancellable, because they highly depend on contextual factors. Conventional implicatures usually resist cancellation because they convey a side comment by the speaker and do not contribute to the at-issue meaning of the sentence (\citealt{Potts2005}, \citealt{SimonsEtAl2010}).  I assume that the neutral meaning of postnominal \textit{cualquiera} as ‘common’ (COM) is the lexical meaning, whereas the derogatory meaning as ‘low value’ (DER) is a contextual or pragmatic interpretation of ‘common’ (i.e. having common properties or properties like anyone or anything else is judged contextually as having unextraordinary properties). The empirical evidence for this assumption is the observation that the derogatory meaning requires contextual information in which some negative attitude is expressed, such as ‘mediocre’. To confirm this assumption, I informally asked several speakers to see which reading could be canceled. The judgments seem to confirm that the derogatory meaning is cancellable. The word \textit{cualquiera} is not translated into English because it was up to the native speakers to decide which meaning best fits the context introduced by the \textit{pero} ‘but’ clause:}

\ea \label{bkm:Ref42976554} %ex 141 
\ea  
\gll \#Juan es un hombre \textit{cualquiera}, pero no es como los demás.\\
{Juan} is a man \textsc{cualquiera} but \textsc{neg} is like the others\\
\glt ‘Juan is a man, but he is not ordinary.’
\ex \label{bkm:Ref42976554-2}
\gll Juan es un hombre \textit{cualquiera}, pero no es un hombre sin nivel.\\
Juan is a man \textsc{cualquiera} but \textsc{neg} is a man without level.\\
\glt{‘Juan is a man, but he is not a man without class.’}
\z
\z

None of the speakers accepted the sentence in \REF{bkm:Ref42976554}. The first clause in \REF{bkm:Ref42976554} expresses that Juan is a man like any other, while the second clause in \REF{bkm:Ref42976554} expresses the opposite, namely that Juan is not like others, resulting in a contradiction. {I assume that it is this contradiction that triggered the negative judgments in \REF{bkm:Ref42976554}.} The sentence in \REF{bkm:Ref42976554-2} does not lead to a contradiction, because it expresses that Juan is like many others, but he is not a man without class or a simple man. I take these results to mean that the meaning ‘common’ is the lexical meaning of \textit{cualquiera} and that the derogatory meaning is a contextual meaning derived pragmatically. {The frequency distribution of the derogatory meaning confirms this assumption because the derogatory meaning is less frequent in the corpus data than the meaning of \textit{cualquiera} as ‘common’ (see \chapref{sec:3.2.1}). However, a controlled experimental study is needed to confirm these results.}  

Assuming that the lexical meaning of \textit{cualquiera} is its free choice meaning that triggers alternatives (see \chapref{sec:2}), a reasonable assumption would be that the meaning ‘common’ (COM) is semantically related to the free choice ‘any’. I will suggest such an analysis in \sectref{sec:3.5}, where I will show that the meaning ‘like anyone else/anything else’ or ‘common’ is a special case of ‘any’. Having common properties is judged contextually as having unremarkable properties.

In the following section, I will investigate further semantic properties of \textsc{un} N \textit{cualquiera} that will help define and characterise its meanings more precisely.

\subsection{{+/–} {evaluative} {\textit{cualquiera}}} \label{sec:3.2.4}

Evaluative \textit{cualquiera} with the readings ‘common’ or ‘normal’ should be distinguished from FC \textit{cualquiera} ‘any’. Both interpretations can be made explicit in the following example:

\ea \label{ex:key:142} %ex 142
\gll la existencia de un delito \textit{cualquiera}\\
the existence of a crime \textsc{cualquiera}\\
\glt ‘the existence of any criminal act’                                     \textit{cualquiera} (FCI)

\glt ‘the existence of a common criminal act’ (EVAL)

  (@ 69)
\z

Under the FCI reading, the probability of any sort of criminal act is calculated as a mean of every single criminal act (including a very infrequent and unusual one such as a murder), whereas the probility of a common criminal act is not the probability of any kind of criminal act but the probability of the occurrence of a certain type of criminal act that can be described as a prototypical, usual, or common criminal act (e.g. robbery). The same observation applies to other readings as well. A criminal act that is similar to many other criminal acts is a certain type of criminal act, namely the one that shares properties with many other criminal acts. An unremarkable criminal act is also a certain type of criminal act. Only in the FCI case can \textit{cualquiera} be modified by a subjunctive relative clause:

\ea \label{ex:key:143} %ex 143
\gll la existencia de un delito \textit{cualquiera} que sea\\
the existence of a crime \textsc{cualquiera} that is.\textsc{SBJV}\\
\glt ‘the existence of any criminal act that there might be’ FCI only

  \#‘the existence of a common criminal act that there might be’

  (@ 70)
\z

Assuming that there are ony three kinds of crimes -- A, B, and C -- and assuming further that A and B are common, whereas C is very infrequent, \textit{un delito cualquiera} under the reading ‘common’ denotes a more restricted set of crimes, namely the set \{A, B\}. FCI \textit{cualquiera} does not add this information about the distribution, whereas \textit{cualquiera} on the ‘common/frequent’ reading does. Set-theoretically, ‘common’ is thus a subset of a set denoted by FCI \textit{cualquiera.} The extension of criminal acts is thus distinct under the FCI and ‘common’ readings, with the ‘common/frequent’ reading denoting a subset of the FCI denotation with respect to its extension.

\ea \label{ex:key:144} %ex 144
The set denotation of \textit{cualquiera} as EVAL is a subset of the FCI ‘any’.
\ea \textit{cualquiera} (FCI) = \{A,B,C\}
\ex \textit{cualquiera} (EVAL ‘COM/EVAL1’) = \{A,B\}
\z
\z

What does \textit{cualquiera} denote by its intension? I suggest analysing the postnominal \textit{cualquiera} in a way similar to adjectives or nominal compounds that modify nouns and describe a property of a property. I call them predicate modifiers (see \sectref{sec:3.3}):

\ea \label{ex:key:145} %ex 145
\gll Juan es un hombre \textit{cualquiera} / ordinario / simplón.\\
Juan is a man \textsc{cualquiera} / ordinary / simple\\
\glt ‘Juan is an ordinary/simple man.’

  (@ 3)
\z

I will show later that this predicate modifier analysis of \textit{cualquiera} can serve us as a basis for explaining some aspects of diatopic variation of Spanish \textit{cualquiera}, in particular some striking differences between European and Argentinian Spanish (see \sectref{sec:3.3} and \chapref{sec:4} for details).

\section{{Postnominal} {\textit{cualquiera}} {as} {a} {predicate} {(modifier)}} \label{sec:3.3}

I assume that postnominal \textit{cualquiera} should be analysed as a postnominal modifier like the postnominal adjective \textit{común}, \textit{ordinario} ‘common, ordinary’:

\ea \label{ex:key:146} %ex 146
 Un hombre [\textsubscript{Modifier} \textit{cualquiera}] = Un hombre [\textsubscript{Modifier} \textit{común/ordinario}].
\z

On this analysis, we should expect that postnominal \textit{cualquiera}, can be coordinated, just like postnominal adjectives. And indeed, we find coordination of \textit{cualquiera} with adjectives such as \textit{mediocre} ‘mediocre’\textit{, ordinario} ‘ordinary’, and \textit{normal} ‘normal’.

\ea \label{ex:key:147} %ex 147
\gll  para un día \textit{cualquiera} y normal en el Ciaam.\\
for a day \textsc{cualquiera} and normal in the Ciaam\\
\glt ‘for a common and normal day in Ciaam’

  (@ 79)
\z

\ea \label{ex:key:148} %ex 148
\gll y no es para llenar un currículo para una hoja de vida ordinaria y \textit{cualquiera}.\\
and \textsc{neg} is for filling a resume for a sheet of life ordinary and \textsc{cualquiera}\\
\glt ‘and it is not to fill in the CV with a page of an ordinary and normal life.’

  (@ 80)
\z

\ea \label{ex:key:149} %ex 149
\gll comprendí que no era más que un ser mediocre y \textit{cualquiera}\\
understood that \textsc{neg} was more than a being mediocre and \textsc{cualquiera}\\
\glt ‘I understood that I was just a mediocre and normal being’

  (@ 81)
\z

However, this kind of coordination is not very frequent in European Spanish (3 total occurrences out of 357 examples), which suggests that the use of postnominal \textit{cualquiera} as an adjective-like category is not very common in this variety of Spanish. The occurrence of definite/quantified nouns with postnominal \textit{cualquiera} as observed in \sectref{sec:3.1} finds a straight-forward explanation in my analysis of \textit{cualquiera}. It is because \textit{cualquiera} acts in a way similar to adjectives like ‘common, unremarkable, worthless’, and adjectives can modify different types of noun phrases, including definite noun phrases, demonstrative and quantified noun phrases.

I will now present a detailed formal analysis of the postnominal \textit{cualquiera} as a predicate (modifier).

\subsection{{Formalisation} {of} {\textit{cualquiera}} {as} {a} {predicate} {(modifier)} } \label{sec:3.3.1}

The idea in a nutshell is that postnominal \textit{cualquiera} as a predicate belongs to different types of predicates, which is why it triggers different interpretations such as ‘common’ (COM) and ‘low value’ or ‘not special’ (DER) observed in \sectref{sec:3.2}. In what follows, I will formalise the different meanings of the predicate \textit{cualquiera}.

Postnominal \textit{cualquiera} with the meaning \textit{común} ‘common’ as ‘usual’ or ‘frequent’ should be analysed as a frequency adjective that modifies nouns with a reference to kinds, as in \REF{bkm:Ref72354905} (see \citealt{Carlson1977}, \citealt{Chierchia1998}, \citealt{GehrkeMcNally2015}). They do not modify proper nouns, as shown in \REF{bkm:Ref72354921}.

\ea \label{bkm:Ref72354905} %ex 151 
Cats are common in Europe.
\z 

\ea \label{bkm:Ref72354921} %ex 152
*Angela Merkel is common in Germany.
\z

The function of a postnominal frequency adjective such as English \textit{common} as in \textit{a} \textit{common example/student} is to turn kinds into subkinds. It takes a kind as an argument and returns a subkind as an output.

At this point, I need to give a short introduction into kind nouns (see \citealt{GehrkeMcNally2015} for this introduction). Nouns denote properties of kinds or individuals, as shown in the representation in \REF{bkm:Ref42889594} (see \citealt{Carlson1977}, \citealt{Chierchia1998}, \citealt{GehrkeMcNally2015}). I use a subscript α\textit{\textsubscript{k}} to indicate a variable that ranges over kinds (see \citealt{GehrkeMcNally2015}). No subscript α indicates a variable that ranges over atomic individuals:

\ea \label{bkm:Ref42889594} %ex 153 
⟦[\textsubscript{NP} [\textsubscript{N} car]]⟧: λ.\textit{x\textsubscript{k}}[car(\textit{x\textsubscript{k}})]
\z

These nouns realise a syntactic Num(ber) projection, which semantically is interpreted as properties of token entities (see \citealt{GehrkeMcNally2015}, and references therein). R is \citegen{Carlson1977} realisation relation. The kind variable is bound off by the existential quantifier in the noun representation in \REF{bkm:Ref42889808}. The numberless noun \textit{car} denotes the set of all kinds of cars (station wagons, sedans, etc.), the (phonologically identical) singular NumP \textit{car} denotes the set of atomic individuals that realise some kind of car, and the plural NumP \textit{cars} denotes the set of token pluralities of some kind of car.

\ea \label{bkm:Ref42889808} %ex 154 
⟦[\textsubscript{NumP} [\textsubscript{NP} car]]⟧: λ\textit{y}${\exists}$\textit{x\textsubscript{k}}[car(\textit{x\textsubscript{k}}) \& R(\textit{y}, \textit{x\textsubscript{k}})]
\z 

I assume that postnominal \textit{cualquiera} can denote a property of kinds, whereas \textit{cualquiera} modifying a noun denotes a property of a property of kinds:

\ea \label{ex:key:155} %ex 155
\ea  ⟦postnominal cualquiera⟧: λ\textit{P}λ\textit{x\textsubscript{k}}[(cualquiera (\textit{P}))(\textit{x\textsubscript{k}})]
\ex ⟦estudiante cualquiera⟧: λ\textit{x\textsubscript{k}}[(cualquiera (student))(\textit{x\textsubscript{k}})]

    \glt ‘kind of a student.’
\z
\z

The experiencer/evaluator function of \textit{cualquiera}, which can be made salient by \textit{para} x ‘for x’ and is represented as a subscript E on the predicate modifier \textit{cualquiera}, which predicates over types of individuals according to some evaluator/experiencer E in a given context C:

\ea \label{ex:key:156} %ex 156
  ⟦estudiante cualquiera⟧\textsuperscript{c} : λ\textit{x\textsubscript{k} }[(cualquiera\textsubscript{E} (student))(\textit{x\textsubscript{k}})]

\glt ‘kind of student for the experiencer/evaluator E’
\z

The following satisfaction condition in \REF{bkm:Ref54641377} expresses that the kind property represented by N \textit{cualquiera} describing a set of elements denoted by the kind is very frequent. Thus, \textit{a common student} denotes a kind of a student such that all students that represent this kind occur very frequently. What it means to be \textit{un estudiante cualquiera} ‘a common student’ is established pragmatically. Thus, \textit{un estudiante cualquiera} denotes a kind of a student that has usual or frequently occurring grades, say \sectref{sec:2.7} (= average grade) in the German system. I adopt the condition of use of the English frequency adjective \textit{common} from \citet{GehrkeMcNally2015} and adapt it for \textit{cualquiera} on the assumption that it has the same use as the Spanish frequency adjectives \textit{común} ‘common’ or \textit{corriente} ‘frequent’. “(Frequency) distribution in the representation is a function that yields the distribution of a set of entities at the given index, with high values” (see \citealt{GehrkeMcNally2015}, my addition of the bracket).

\ea \label{bkm:Ref54641377} %ex 157
Condition of use in a context for \textit{cualquiera} with the meaning ‘common/frequent’:

    ${\forall}$\textit{P}, \textit{x\textsubscript{k}}, \textit{i}[(cualquiera(\textit{P}))(\textit{x\textsubscript{k}}) at \textit{i} $\leftrightarrow $ [\textit{P}(\textit{x\textsubscript{k}}) \& distribution(\{\textit{y} : R(\textit{y}, \textit{x\textsubscript{k}}) at \textit{i}\}) = \textit{high}]]
\z 

\ea \label{ex:key:158} %ex 158 
\ea ⟦estudiante cualquiera⟧: λ\textit{x\textsubscript{k}} [(cualquiera({student}))(\textit{x\textsubscript{k}})]
\ex ${\forall}$\textit{x\textsubscript{k}}, \textit{i}[(cualquiera(student))(\textit{x\textsubscript{k}}) at \textit{i} $\leftrightarrow $ [student (\textit{x\textsubscript{k}}) \& distribution(\{\textit{y} : R( \textit{y}, \textit{xk}) at \textit{i}\}) = \textit{high}]]
\z
\z 

The type N \textit{cualquiera} implies a contrast to another type, which denotes a singleton and has a low (frequency) distribution. The satisfaction condition for the meaning N non\textit{{}-cualquiera} or N \textit{distinguido/especial} ‘special/particular N’ are the same as for N \textit{cualquiera} besides the low frequency:

\ea \label{ex:key:159} %ex 159
    ${\forall}$\textit{P}, \textit{x\textsubscript{k}}, \textit{i}[(distinguido(\textit{P}))(\textit{x\textsubscript{k}}) at \textit{i} $\leftrightarrow $

    [\textit{P}(\textit{x\textsubscript{k}}) \& distribution(\{\textit{y} : R( \textit{y}, \textit{x\textsubscript{k}}) at \textit{i}\}) = \textit{low}]]
\z 

\ea \label{ex:key:160} %ex 160
\ea ⟦estudiante distinguido⟧: λ\textit{x\textsubscript{k}} [(distinguished({student}))(\textit{x\textsubscript{k}})]
\ex ${\forall}$\textit{x\textsubscript{k}}, \textit{i}[(distinguished(student))(\textit{x\textsubscript{k}}) at \textit{i} $\leftrightarrow $ [student (\textit{x\textsubscript{k}}) \& distribution(\{\textit{y} : R( \textit{y}, \textit{xk}) at \textit{i}\}) = \textit{low}]]
\z
\z

Under this approach, the contribution of \textit{cualquiera} is to restrict the kind denoted by the noun N, that is, to turn a kind into a subkind. For instance, the noun \textit{estudiante} denotes a kind of a student and the NP \textit{estudiante cualquiera} denotes a common kind of student as opposed to a non-common one. On the extension, N denotes a set of students that belong to the kind ‘student’ and \textit{estudiante cualquiera} denotes a restricted set of students, namely a set that is described as a ‘common student’, rather than as a ‘distinguished student’.

In case N denotes an event noun (indicated by the subscript α\textit{\textsubscript{e}}), \textit{cualquiera} denotes a property of a property of an event (see \citealt{GehrkeMcNally2015} for predicate modifiers of event nouns):

\ea \label{ex:key:161} %ex 161
\ea  ⟦cualquiera⟧: λ\textit{P}λ\textit{x\textsubscript{e}}[(cualquiera(\textit{P}))(\textit{x\textsubscript{e}})]
\ex ⟦evento/historia cualquiera⟧: λ\textit{x\textsubscript{e}}[(cualquiera(event/story))(\textit{x\textsubscript{e}})]

    \glt ‘frequent event/story, a story/event that frequently occurs’

  Satisfaction condition:
\ex ${\forall}$\textit{x\textsubscript{e}}, \textit{i}[(cualquiera(story))(\textit{x\textsubscript{e}}) at \textit{i} $\leftrightarrow $ [story(\textit{x\textsubscript{e}}) \& distribution(\{\textit{y} : R( \textit{y}, \textit{x\textsubscript{e}}) at \textit{i}\}) = \textit{high}]]
\z
\z

To sum up, \textit{cualquiera} in DET [+/– def.] N \textit{cualquiera}, which can be coordinated with adjectives, is not a determiner, but a nominal predicate that describes the distribution of properties among animate entities (non-temporal frequency) or the distribution of properties of events in time (temporal frequency with the meaning \textit{common} and \textit{frequent in time}).

On this analysis, N \textit{cualquiera} in its COM meaning can be used with determiners and quantifiers as is expected from an adjectival use of \textit{cualquiera} (see \REF{bkm:Ref58144031}, \REF{bkm:Ref63546096}, and \REF{bkm:Ref63545986}–\REF{bkm:Ref63545989} for data):

\ea \label{ex:key:162} %ex 162
\ea  ⟦ese/el estudiante cualquiera⟧: ι \textit{x}[(cualquiera(student))(\textit{x})]
\ex ⟦algún/otro/un estudiante cualquiera⟧: some \textit{x}[(cualquiera(student))(\textit{x})]
\ex ⟦mucha gente cualquiera⟧: many \textit{x}[(cualquiera(people))(\textit{x})]
\z
\z

As \textit{cualquiera} modifies an NP, which denotes a kind of entity or individual, the predicative N \textit{cualquiera} can refer to some specific type/kind of individual, such as the common type. This is indeed what we find in the data: \textit{el hombre cualquiera, la mujer cualquiera} ‘the common type of men or women’ as opposed to \textit{el hombre distinguido/particular} ‘the distinguished type of men’ (see 3.3.1).

In the following example, repeated from \REF{bkm:Ref63545986}, the definite noun refers to a certain kind of man, namely to the common man {as a type}. One could also describe this type as \textit{el hombre de la calle} lit. ‘the man from the street’ or ‘common/normal man’:

\ea \label{bkm:Ref63545986} %ex 163
El tipo de el C1, que antes fue el de el [sic] C2, es un timo. Es el prototipo de pícaro español. Un personaje que nuestros más célebres escritores de el [sic] Siglo de Oro describieron con crudo realismo – cómo olvidar el Lazarillo de Lázaro de Tormes –.\\ 

\gll {El} {hombre} \textit{cualquiera} {del} {autobús} {no} {siente} {vergüenza} {alguna,} {aunque} {debería,} {por} {tratar} {de} {engañar} {a} {otras} {personas} {que} {están} {dispuestas} {a} {ayudar} {a} {aquellos} {extraños} {que} {lo} {pasan} {mal}.\\
the man \textsc{cualquiera} of.the bus \textsc{neg} feels shame any although should for trying to deceive to other people that are willing to help to those strangers that it pass badly\\
\glt ‘[…] The common man from the bus does not feel any shame, although he should,…’

  (@ 82)
\z

The same analysis applies to \textit{el día cualquiera}, namely the most prototypical day of the week:

\ea \label{ex:key:164} %ex 164
\gll diría que es el día más indeterminado, el día \textit{cualquiera} de la semana, un día metido en la selva de los días\\
would.say that is the day most indeterminate, the day \textsc{cualquiera} of the week, a day stuck in the jungle of the days\\
\glt ‘I would say that it is the most undistinguished day, just some day of the week, a day put in the jungle of the days.’

  (@ 100)
\z

The definite article can thus refer to a specific type of individual, namely a common one, and not some other type, a non-common one.

We also find demonstratives, which refer to a specific type of a man of speaker’s reference that has prototypical properties of a man. The example is embedded in a philosophical context, which contains a statement about a certain type of a man:

\ea \label{bkm:Ref63545989} %ex 165
\gll Las armas de este hombre \textit{cualquiera}, de aspecto normal, se resumen en: prudencia\\
the weapons of this man \textsc{cualquiera}, of appearance normal, \textsc{refl} summarised in: prudence.\\
\glt ‘The weapons of this common, normal-looking man are summarised in: caution.’

  (@ 83)
\z

Another consequence of the kind analysis of postnominal \textit{cualquiera} is that it explains the lack of occurence with degree modifiers or adverbs like \textit{muy} ‘very’ as in \textit{un hombre *(muy) cualquiera} ‘a common man’. Degree adverbs like \textit{muy} ‘very’ combine with gradable adjectives, not with kind predicates \citep{Kennedy1999}.

To recap, postnominal \textit{cualquiera} has been analysed as a frequency adjective ‘common’ in line with \citet{GehrkeMcNally2015}. Its lexical meaning is to describe kinds of individuals or entities that have a high frequency, that is, considering all invidividuals or entities denoted by N, we would observe very often a set of individuals or entities that are described by N \textit{cualquiera}.

\subsection{{Formalisation} {of} {the} {derogatory} {meaning} {of} {N} {\textit{cualquiera}}} \label{sec:3.3.2}

The idea is that the derogatory meaning of postnominal \textit{cualquiera} has an additional satisfaction condition on top of describing individuals or types of individuals with a high-frequency distribution. Namely, the second satisfaction condition in \REF{bkm:Ref58094203} expresses that N \textit{cualquiera} in its DER reading, such as \textit{equipo cualquiera} ‘ordinary team’ denotes a type of individuals that is ranked lower on some qualitative scale according to the experiencer E expressed in \textit{para} x ‘for x’ than a team that is not described as \textit{cualquiera} or a distinguished team (e.g. \textit{equipo con Messi} ‘team with Messi’):

\ea \label{bkm:Ref58094203} %ex 166
⟦equipo cualquiera\textsubscript{derogatory} ⟧\textsuperscript{c}: λ\textit{x\textsubscript{k}}[(cualquiera\textsubscript{E} (\textit{team}))(\textit{x\textsubscript{k}})]
\ea  ${\forall}$\textit{x\textsubscript{k}}, \textit{i}[(cualquiera(team))(\textit{x\textsubscript{k}}) at \textit{i} $\leftrightarrow $ [team(\textit{x\textsubscript{k}}) \& distribution(\{\textit{y} : R( \textit{y}, \textit{xk}) at \textit{i}\}) = \textit{high}]]
\ex \label{bkm:Ref58094203b} ${\forall}$y\textsubscript{k,} ${\forall}$x\textsubscript{k,} \textit{i}[(cualquiera′ (y\textsubscript{k}) \& distinguished′(y\textsubscript{k}) at \textit{i} $\leftrightarrow $ [team(\textit{x\textsubscript{k}})(\textit{y\textsubscript{k}}) \& (\textit{x\textsubscript{k}}) <\textsubscript{lower on some pragmatic scale according to E} (\textit{y\textsubscript{k}}) at \textit{i}]]
\z
\z

On this analysis, \textsc{un} N \textit{cualquiera} cannot have an evaluative meaning in contexts that do not include any evaluation on a pragmatic scale according to an experiencer, such as mathematical contexts. This prediction is borne out empirically:

\ea \label{ex:key:167} %ex 167
  \gll Es un cuerpo geométrico \textit{cualquiera}. \# Sin nivel.\\
is a body geometric \textsc{cualquiera}. \# without level.\\
  \glt  ‘It’s an unremarkable geometric body. \# Without class.’

  (@ 101)
\z

I will show in \chapref{sec:4} that in some languages or language varieties DER is only felicitous under the condition \REF{bkm:Ref58094203b}; that is, the requirement that the set of individuals or entities must be described as usual/frequent withouth the condition b is not given. French \textit{n’importe quoi} (see \ref{bkm:Ref72355051-1}) and Argentinian Spanish \textit{cualquiera} (see \ref{bkm:Ref72355051-2}) are such a case. They necessarily have a derogatory meaning. According to the derogatory meaning of \textit{n’importe quoi}/\textit{cualquiera,} the evaluator considers that saying something undistinguished is worse than saying something distinguished (see \chapref{sec:4} for a detailed analysis):


\ea \label{bkm:Ref72355051} %ex 168
\ea{ \label{bkm:Ref72355051-1}
\gll Marie a dit (du grand) \textit{n’importe quoi}.  \\
Marie has said (of.the big) anything.\\}\jambox*{(French)}
    \glt ‘Marie said a big bullshit.’

      (@ 30)
\ex{\label{bkm:Ref72355051-2}
\gll Me mandaste \textit{cualquiera}.  \\
to.me sent \textsc{cualquiera}\\}\jambox*{(ArgSp)}
    \glt ‘You told me bullshit.’

    (@ 102)
\z
\z

\subsection{{The} {FCI} {analysis} {of} {postnominal} {\textit{cualquiera}} {in} {predicative} {contexts}} \label{sec:3.3.3}

The goal of this section is to adapt the FCI analysis of postnominal \textit{cualquiera} in predicative contexts as in \textit{es un día cualquiera} and to show the connection between FCI ‘anyone’ or ‘anything’ and the evaluative reading ‘common’ (COM).

  Consider the following two examples:

\ea \label{bkm:Ref72353992} %ex 169 
\gll Juan es un estudiante \textit{cualquiera} de mi clase.\\
Juan is a student \textsc{cualquiera} of my class.\\
\glt ‘Juan is an ordinary student in my class.’

  (@ 74)
\z


\ea \label{bkm:Ref72354006}  %ex 170
\gll Juan es un estudiante como \textit{cualquiera} de mi clase.\\
Juan is a student like \textsc{cualquiera} of my class.\\
\glt ‘Juan is a student like everyone in my class.’

  (@ 75)
\z

One possible analysis of the two aformentioned examples is to interpret them as semantically equivalent. Comparative constructions are known to license FCIs (see \citealt{Aloni2021} and \sectref{sec:2.2}). In this view, \textit{un hombre cualquiera} would be interpreted as ‘a man like anyone [=any man]’. The function of postnominal \textit{cualquiera} in \REF{bkm:Ref72353992} is thus to express a predicate that has a similar semantic interpretation as the pronoun \textit{cualquiera} in an overt comparative construction as in  \REF{bkm:Ref72354006} (for comparative clause analyses of cases like in \REF{bkm:Ref72354006}, see \citealt{Bresnan1973}, \citealt{Napoli1986}, \citealt{Saez1992}, \citealt{Kennedy1999}, among many others).

Different analyses have been suggested in the literature on comparative clauses (see \citealt{Bresnan1973}, \citealt{Napoli1986}, \citealt{Saez1992}, \citealt{Kennedy1999}, among many others). I will follow the coordination hypothesis of comparatives suggested by \citet{Napoli1986} and \citet{Saez1992}, among others, who assume that comparatives represent coordinations (on either the phrasal or sentence level); that is, \textit{como cualquiera de mi clase} is represented as a coordinated phrase or clause (\cite{cp2009}). I assume that the comparative phrase in \REF{bkm:Ref72354006} is clausal because it can be paraphrased as a clause. In order to represent the clausal analysis, I mark the predicate as deleted as shown in the English translation (see \citealt{Bresnan1973}):

\ea \label{bkm:Ref72354047}  %ex 171

\gll Juan es un estudiante como lo es \textit{cualquiera} de mi clase.\\
  Juan is a student like it is \textsc{cualquiera} of my class\\
\glt ‘Juan is a student like everyone in my class \sout{is a student}.’

  (@ 76)
\z

Let me suggest a simple syntactic analysis of \REF{bkm:Ref72354047}. I analyse \textit{cualquiera de mi clase} as a subject of a copulative clause, which is inside the comparison clause introduced by \textit{como}. Assuming that \textit{como} conjoins two sentences (see \citealt{Napoli1986}, \citealt{Saez1992}), \REF{bkm:Ref72354006} can be interpreted as a coordination of two conjoined copulative sentences on the semantic level:

\ea \label{ex:key:172} %ex 172
 ‘[Juan is a student] and [[\textsubscript{subj.} everyone in my class] is a student].’
\z

The syntax of \textit{cualquiera} that expresses a comparison is that of a {comparison clause}. This comparison clause is conjoined with the main clause (TP) due to the presence of the comparative particle \textit{como}. In \REF{bkm:Ref72354100}, the first clause has \textit{Juan} as a subject and a predicative DP \textit{un estudiante.} The comparative clause contains \textit{cualquiera} as a subject and a clitic pronoun \textit{lo} which is coreferent with the predicative complement inside the matrix clause:

\ea \label{bkm:Ref72354100} %ex 173
[\textsubscript{TP} Juan es [\textsubscript{DP} un estudiante]\textsubscript{j}] [\textsubscript{\& P} como lo\textsubscript{j} es cualquiera de mi clase]

\glt ‘Juan is a student, and everyone (else) is a student in my class.’
\z

\largerpage
\REF{bkm:Ref72354100} expresses similarity with respect to some (prototypical) property P* of a student, such as with respect to Juan’s qualifications or grades, which can be made salient in the context as will be demonstrated in more detail in \REF{bkm:Ref54682554}. In a less strict interpretation, the sentence in \REF{bkm:Ref72354100} with the assumption that the comparative clause represents a conjoined sentence as suggested here could simply mean that Juan and everyone in the speaker’s class are students or have the property of being a student (as opposed to teachers, for example). But this interpretation is too weak and makes false predictions for the use of sentences such as \REF{bkm:Ref72354100} in the following context. Imagine a scenario where the speaker wants to say that Juan is a certain type of a student, namely an average student. In this scenario, the meaning ‘Juan is a student and everyone else in my class is a student’ is underinformative. The solution to this problem is pragmatics. What it means to be a student and to be similar to other students is established pragmatically by some property P*; that is, the comparison of students is a comparison with respect to some (prototypical) property P* that holds for students. Thus, the sentence in  \REF{bkm:Ref72354100} states that Juan and other students in the class of speaker’s reference have the same prototypical property P* as students, which can be usual grades, for instance. From this analysis, we can already pragmatically infer that Juan is an average or unremarkable student. If Juan were a brilliant student, he probably would not be like every other student in the speaker’s class, because it is very unlikely that the whole class consists of brilliant students.

I will now formalise the semantic analysis of postnominal \textit{cualquiera} as the pronoun \textit{cualquiera} in a comparative clause analysis.

Given the syntax in \REF{bkm:Ref72354100}, the semantic representation entails a universal quantifier \textit{cualquiera} inside the comparative clause \& P. In this case, \textit{cualquiera} is interpreted as a universal quantifier with the meaning ‘everyone’, ranging over individuals in the speaker’s class, and these individuals represent a comparison class of students to whom the subject of the main clause (here: Juan) is compared:

\ea \label{bkm:Ref54682554} %ex 174
\textit{cualquiera} in a comparison clause
\ea Juan es un estudiante [\textsubscript{ModifP} \textit{cualquiera}]. Tiene calificaciones normales.\\ 
‘Juan is a student like any other. He has normal grades.’
\ex Semantic interpretation: Juan has property P*$\land$ $\forall$x [x $\in$ my class $\rightarrow$ x has property P*]
%Juan is a student with a property P* and [${\forall}$]([anyone in my class, λx. has it (= the property P*)])
\ex Prose of b: Juan is a student with property P* and for each student x in my class, x has \textit{it} [=P*].\footnote{{Another
    way to interpret the conjoined clause with} {\textit{cualquiera}} {as a subject is to assume that it implies a covert modal or attitude verb that links the expression of similarity to the speaker’s evaluation. In this case, we would have universal quantification over a modal verb that describes speaker’s attitude (see \sectref{sec:2.2.2}).} {\textit{Juan es un hombre cualquiera} ‘Juan is a man \textsc{cualquiera}’} {would be interpreted similar to} {\textit{Juan es un hombre cualquiera para mí} ‘According to me, Juan is a man \textsc{cualquiera}.’}
    
  \ea
  {It seems to me that Juan is a man with P* and for each student, he/she has P* too.}
  \z
}
  \z
\z

The positive outcome of a comparison clause analysis is that it explains the licensing of postnominal \textit{cualquiera} in affirmative predicative contexts. Recall from the puzzle described in \sectref{sec:2.3} that affirmative predicative and episodic contexts were considered problematic for FCIs such as \textit{cualquiera}. However, in my analysis, the licensing of postnominal \textit{cualquiera} under indicative predicative verbs is explained by its predicate that refers to similarity with all other students from my class. Recall that FCIs are generally licensed in universal contexts in Spanish and other languages (see \sectref{sec:2.1} and \sectref{sec:2.2}).

Moreover, the suggested analysis explains why the ‘common’ reading cannot be cancelled (see \sectref{sec:3.2.3}). The ‘common’ reading is the result of analysing postnominal \textit{cualquiera} as a modifier that contains a universal quantifier that licenses \textit{cualquiera} in predicative contexts. 

The suggested analysis also explains why N \textit{cualquiera} is not gradable in European Spanish, as shown in \REF{bkm:Ref58074686}. Postnominal \textit{cualquiera} is a predicate with quantificational force of the same type as ‘every(one), any(one)’ and not a gradable adjective like ‘ordinary’, which is why it cannot be modified by degree adverbs like \textit{muy} ‘very’ (see \chapref{sec:2}):

\ea \label{bkm:Ref58074686} %ex 175
\gll Es un estudiante (*muy) \textit{cualquiera}.\\
is a student (very) \textsc{cualquiera}\\
\glt ‘He is a student (*very) like anyone else.’

  (@ 77)
\z

\largerpage
Another prediction of the comparative clause analysis is that it should be possible to modify postnominal \textit{cualquiera} by adverbs like \textit{casi} ‘almost’ that restrict the domain of quantification of \textit{cualquiera de mi clase} to ‘almost everyone in my class’ (see \sectref{sec:2.1} for such adverbial modification of \textit{cualquiera}). However, the judgements from native speakers of European Spanish and the absence of \textit{un N casi cualquiera} in the corpus data in European Spanish (despite some occurrences in Latin American Spanish), seem to suggest that the postnominal \textit{cualquiera} is not modifiable by the modifier \textit{casi} in European Spanish. However, this result needs to be tested in a large experimental study in future research.

\ea \label{ex:key:176} %ex 176
  \gll Fue un día *(casi) \textit{cualquiera}.\\
  was a day (almost) \textsc{cualquiera}\\
  \glt \#‘It was a day like almost any other day.’\footnote{{There
        are other cases of adverbial modification of the free choice item} {\textit{cualquiera}}{, which shows that} {\textit{cualquiera}} {is a quantifier with the meaning ‘any’. However,} {\textit{cualquiera}} {appears in these examples prenominally as a determiner and not as a postnominal modifier (see \sectref{sec:4} for further examples of the determiner} {\textit{cualquiera}} {in predicative position):}

    \ea \gll No tiene nada de particular; pero por eso mismo parece Tamayo, porque no tiene nada de particular. \textit{Parece} \textit{cualquier} \textit{cosa} \textit{menos} \textit{un} \textit{gran} \textit{poeta.}\\
    \textsc{neg} has nothing of particular; but for that same seems Tamayo, because \textsc{neg} has nothing of particular seems \textsc{cualquier} thing except a great poet\\
    \glt {‘He does not have any particular properties. But for the same reason he resembles Tamayo, because he does not have any particular property. He resembles} {\textit{anything but a great poet}}{.’} 
    {(Solos} {de} {Clarín}, {Leopoldo} {Alas)} 
    \z 
    }\footnote{According to two native speakers of European Spanish, acceptable according to a Colombian Instagram user and a Mexican Spanish writer Edmée Pardo ({\url{https://www.instagram.com/piaguzmanv/}{ {and}} \url{https://edmeepardo.com/bibliomemoria-fichas/rfn6txw4j6pg3l2-6kdbj-2e74j-d9r63-zmw5n-az8dg-4e5j4-bamwc-d3fda-c6ygp-efe6w-nr8wf-tnmwk-mwjcd-4wp56}) {} and also acceptable for Michael Robinson, see \textit{Prólogo} in Pérez Gellida, \textit{Memento mori}}.}
\z

The lack of frequent modification by \textit{casi} in \REF{ex:key:176} suggests that the postnominal modifier \textit{cualquiera} is not syntactically a free choice pronoun like the pronoun \textit{cualquiera}, which is modifiable by \textit{casi} very frequently. Consequently, the postnominal \textit{cualquiera} needs to be analysed as a predicate modifier that is semantically related to the pronoun \textit{cualquiera} in the description of the property as ‘like anyone else’.

The following example shows that once the postnominal \textit{cualquiera} cooccurs with an overt comparative clause like \textit{como cualquier mañana} ‘like any other morning’, the interpretation of postnominal \textit{cualquiera} changes into the meaning ‘ordinary’ or ‘just normal’, which corresponds to the derogatory reading DER. This suggests that the expression of similarity between the actual referent and a comparison class is not part of the lexical semantic meaning as suggested in the discussion of \REF{bkm:Ref42976554}, because otherwise the following sentence would express the same information twice and it would be redundant.

\ea \label{ex:key:177} %ex 177
\gll Era una mañana \textit{cualquiera}, como \textit{cualquier} mañana.\\
was a morning \textsc{cualquiera} like \textsc{cualquier} morning\\
\glt ‘It was an ordinary morning, like any other morning.’\footnote{{I thank Josep Quer for this data point.} }
\z

The issue of what kind of status the similarity has (lexical semantic or pragmatic implicature) (\citealt{Aloni2021}, \citealt{am2018}) needs to be analysed in future research.

\section{{Summary} {and} {discussion} {of} {the} {predicate} {modifier} {analysis}} \label{sec:3.4}

In \sectref{sec:3.3}, I suggested a predicate modifier analysis of postnominal \textit{cualquiera} with the meaning COM that can explain a certain type of data such as [Determiner/Quantifier N \textit{cualquiera}] and coordination with adjectives. There are several issues that need to be addressed here.

First, recall from \sectref{sec:2.1.3} that there is the standard assumption in the literature that postnominal \textit{cualquiera} occurs only with indefinite nouns and is hence accompanied by the indefinite article, but not by strong determiners (see, in particular, \citealt{Lapesa1964}, \citealt{Rivero2011}). Linguists have suggested different explanations for this restriction. According to one explanation, postnominal \textit{cualquiera} is itself a strong (quantificational) determiner, and this is why it is incompatible with other strong determiners \citep{Lapesa1964}. Another semantic explanation is that N \textit{cualquiera} has an anti-singleton restriction (see \citealt{Menéndez-Benito2010}), making it incompatible with definite determiners or demonstratives, as these elements presuppose a uniquely identified referent (singleton restriction) and thus do not co-occur with \textit{cualquiera} as an FCI. However, contrary to these preddictions, I was able to find examples in which postnominal \textit{cualquiera} is preceded by definite determiners, and which my predicate modifier analysis is able to explain.

Second, there are indeed a lot more data with \textsc{un} N \textit{cualquiera} than with D\textsubscript{def} N \textit{cualquiera} in the corpus (6,938 indefinite occurrences vs. 171 definite occ.). A simple explanation for this preference under copular verbs is that N \textit{cualquiera} is used predicatively. It is a common observation in the literature that predicative NPs are expressed by NPs that are indefinite rather than definite (see \cite{Geenhoven2005}, among others):

\ea \label{ex:key:178} %ex 178
 \gll Era un / *el día normal (para mí).\\
was a / the day normal (for me)\\
\glt ‘It was a/*the usual day (for me).’

  (@ 103)
\z

Third, my analysis predicts that it should be possible to find cross-linguistic evidence for quantifiers reanalysed as nominal predicates. We actually find few cases in Italian and German, where certain quantifiers can be used as nominal predicates in predicative position similar to \textit{un/una cualquiera} in Spanish (see \sectref{sec:2.1}). The only difference to \textsc{un} N \textit{cualquiera} is that the quantifiers in [indefinite + quantifier], such as in \textit{un cualquiera} in Spanish, \textit{un} \textit{niente} in Italian, and \textit{ein Jedermann} in German, have the syntactic status of a noun rather than of a predicate modifier. However, the semantic analysis is the same; that is, the quantifier is interpreted not as a quantifier but as a noun with a property interpretation that describes certain kinds of people. \textit{Ein Jedermann} is a kind of man, as is \textit{un niente} and \textit{un cualquiera.}

\ea{ \label{ex:key:179} %ex 179
\gll Gianni è un niente.  \\
Gianni is a nothing\\}\jambox*{(Italian)}
\glt ‘Gianni is a nobody.’

  (@ 104)
\z 

\ea{ \label{ex:key:180} %ex 180
 \gll Ein Jedermann.   \\
a everyone\\}\jambox*{(Title of a book by an Austrian author)}
\glt ‘Everyman.’
\z

Fourth, my analysis allows us to better approach some aspects of the diatopic variation of \textit{cualquiera} in Spanish. The element \textit{cualquiera} in Argentinian Spanish can appear as a (degree) predicate over nouns that refer to atomic individuals such as \textit{Juan} or that refer, for instance, to a certain movie (see \sectref{sec:2.4}):

\ea{ \label{ex:key:181} %ex 181
\gll La peli /Juan es (muy) \textit{cualquiera}.   \\
the movie /Juan is (very) \textsc{cualquiera}\\}\jambox{(ArgSp/*EurSp)}
\glt ‘This movie/Juan is (very) ordinary/bad.’

  (@ 5 / @ 105)
\z

The difference between Argentinian Spanish and its European counterpart is that \textit{cualquiera} is used as a predicate modifier in European Spanish, but not as a predicate (i.e. a predicative adjective). I will explain this difference by the assumption that the analysis of \textit{cualquiera} as an adjective in Argentinian Spanish is at a more advanced stage than in European Spanish (for details, see \chapref{sec:4}). Syntactically, it can form a regular copula structure, which can be sketched as follows:

\ea \label{ex:key:182} %ex 182
  [\textsubscript{TP} Juan\textsubscript{j} T° es [\textsubscript{PrP} t\textsubscript{j} [\textsubscript{PrP'} Pr° cualquiera]]].

\glt ‘Juan is ordinary/unremarkable.’
\z

I assume that in European Spanish \textit{cualquiera} is not used as a predicate that describes properties of individuals. Let me btriefly outline the syntax of \textit{cualquiera} in its COM reading in European Spanish. In European Spanish, postnominal \textit{cualquiera} is a nominal modifier similar to adjectives like ‘common/unremarkable’. I use the labels Modif(ier) for all kinds of modifiers (adjectives, relative clauses, etc.). This type of postnominal \textit{cualquiera} modifies nouns that denote kinds and, as such, N \textit{cualquiera} can be represented as a KindP (see \citealt{Zamparelli2000} for KindPs):

\ea \label{ex:key:183} %ex 183
  [\textsubscript{DP} un [\textsubscript{KindP} \textsubscript{Kind°} hombre [\textsubscript{ModifP} \textit{cualquiera/cumún/ordinario} ]]]

\glt ‘a common/ordinary man’
\z

The whole DP can appear as a predicative complement of the copula verb \textit{be} in European Spanish, but it cannot appear as a predicative complement by itself as in Argentinian Spanish as will be shown in \sectref{sec:4}): 

\ea{} \label{ex:key:184} %ex 184
  [\textsubscript{TP} Juan\textsubscript{j} [\textsubscript{T'} [T° es [\textsubscript{PrP} t\textsubscript{j} [\textsubscript{PrP'} Pr°\textsubscript{} [\textsubscript{DP} un [\textsubscript{KindP} hombre [\textsubscript{ModifP} \textit{cualquiera/cumún/ordinario}]]]]]]]]

\glt ‘Juan is an ordinary man.’
\z

It can also appear as a nominal modifier of an NP that is the subject of a clause:

\ea{} \label{ex:key:185} %ex 185
  [\textsubscript{DP} \textit{El} [\textsubscript{KindP} \textit{hombre} [\textsubscript{ModifP} \textit{cualquiera}]]] del autobús no siente vergüenza alguna.

\glt ‘The common man from the bus does not feel any shame.’
\z

As we will see in \chapref{sec:4}, the Argentinian Spanish \textit{cualquiera} is not restricted to the use of a nominal modifier as in European Spanish.

\section{{Chapter} {summary} {and} {outlook} {for} {subsequent} {chapters}} \label{sec:3.5}
\largerpage
I started with different observations about \textsc{un} N \textit{cualquiera} in the corpus (see \sectref{sec:3.1}–\ref{sec:3.4}) and showed that in affirmative predicative contexts under indicative verbal mood, it has usually an evaluative meaning that manifests itself in various readings, namely ‘similar to anyone/anything else’ or ‘common’ (COM), and ‘of low value’ (DER) (\sectref{sec:3.2}).


I showed in \sectref{sec:3.3} that ‘similar to anyone/anything else’ or ‘common’ (COM) is a special case of interpreting \textit{cualquiera} as a predicate with quantificiational force similar to the quantifier pronoun \textit{cualquiera} that occurs in a comparison clause and that the DER reading can be derived pragmatically from the neutral reading of ‘common’ or ‘similar to anyone/anything else’.

In the following chapters, I will examine the diatopic and diachronic variation of \textit{cualquiera} in more detail. \chapref{sec:4} will deal with diatopic variation. Chapters \ref{sec:5} will deal with diachronic variation.

The question that will guide my analysis in subsequent sections is why a generalised quantifier with the free choice meaning ‘any’ changes into a predicate modifier as demonstrated in \sectref{sec:3} for European Spanish or a gradable predicate as demonstrated for Argentinian Spanish. My short answer to this question is that it is the predicative context, the negation, and the adjectival appositives that trigger the reanalysis of postnominal \textit{cualquiera} as an FCI into a predicate modifier or predicate of an adjective type. Adjectival appositives can elicit different readings of FCIs in predicative position such as ‘similar to everybody else’, ‘common’, and ‘of low value’, and by doing this, they can turn generalised quantifiers of type <{}<e,t>,t> into predicate modifiers of type <{}<et>,<e,t>{}> and can thus turn the interpretation of ‘any’ into the adjectival interpretation of ‘any’ or ‘common/frequent/normal’. This is how postnominal \textit{cualquiera} gets its adjectival or kind flavour. The adjectival interpretation can be contextually elicited by appositives that specify the property as ‘average/frequent’. The derogatory flavour of ‘nothing special’ is just a negative interpretation of ‘any’ or of ‘average’.

It is a common observation in the literature that semantic change is often triggered by pragmatic enrichment (see \sectref{sec:2.7} and \citealt{Traugott2002}, \citealt{Eckardt2006}, among others). In the case of \textit{cualquiera}, the most crucial trigger for semantic change is syntax, in particular the predicative context and negation, which triggers a contrastive reading between the property described by N \textit{cualquiera} and another contrasting property like N\textit{extraordinario}. The function of adjectival appositives that describe the property of postnominal \textit{cualquiera} as ‘common, frequent, low value’ elicits the adjectival interpretation. The diachronic path of \textit{cualquiera} can be sketched as follows:

\ea \label{ex:key:186} %ex 186
  Diachronic development of \textit{cualquiera}
\begin{enumerate}
\item  Use of \textit{cualquiera} as FCI in modal constructions.
\item   Use of \textit{cualquiera} in predicative position (under negation) with the reading ‘common' or the derogatory reading.
\item   Reanalysis of FCI \textit{cualquiera} as a property of kind nouns in predicative contexts.
\item  Reanalysis of \textit{cualquiera} as a gradable property of proper nouns.
\end{enumerate}
\z

I argued that a small set of data of postnominal \textit{cualquiera} in European Spanish shows that it has reached step 3. In \chapref{sec:4}, I will show that Argentinian Spanish has reached step 4. In Chapter \ref{sec:5}, I will look at the diachronic conditions that have enabled step 3 in European Spanish, namely the change of \textit{cualquiera} from a relative clause into a compound which at some point was possible not only under modal verbs, but also under affirmative predicative verbs.

In the following chapters, I will study the diatopic and diachronic variation in more detail. According to my analysis, \textit{cualquiera} as FCI should be diachronically prior to readings in affirmative predicative sentences. The evaluative readings in some varieties of Spanish or other Romance or Non-Romance languages are expected to occur later than the reading ‘any’ (FCI). Moreover, I will also look into the development of EVAL ‘nonsense’ in Argentinian Spanish \textit{cualquiera} in \chapref{sec:4}.

